\section{Cholesky：利用正定结构}
\label{sec:08-Numerical-Analysis/05-Numerical-Linear-Algebra/03}

一段 Gaussian process 的超参数优化循环。每一轮给定长度尺度与噪声方差，按核函数组装两千阶的核矩阵 \(K\)，然后算对数边际似然，需要的是 \(K\alpha=y\) 的解和 \(\log\det K\)。这个循环要跑几百轮，所以第一件事是别用通用求解器。第二件事是循环跑到第三十七轮时抛出了异常：矩阵不是正定的。

最省事的处理是在对角上加 \(10^{-6}\)，让分解通过，循环继续，优化收敛。这一节想说明的是，那次失败本来是整段代码里信息量最大的一个事件；而在它之前，选择 Cholesky 这件事本身已经换回了三样东西。

\subsection{一个结构假设换回三样东西}

若实对称矩阵 \(A\) 对一切 \(x\ne0\) 满足 \(x^{\trans}Ax>0\)，则存在唯一的下三角矩阵 \(L\)，对角元全为正，使得

\[
A=LL^{\trans}.
\]

第一样是省。对称性意味着上三角是下三角的镜像，一半的乘加根本不必做，一半的存储也不必留。运算量从一般 LU 的 \(\tfrac23n^3\) 降到 \(\tfrac13n^3\)，只存一个三角因子。

第二样是不必选主元。上一节的部分选主元是为了控制乘子的量级，代价是每列扫一遍、交换行、打乱按列分块的访问顺序。正定性把这项开销整个免掉：每一步开方的那个量自动为正，对角元天然就是合格的主元。这里还有一层：交换行会破坏对称性，而对称性正是省下那一半运算的依据，所以选主元与结构复用本来是两难的，正定性让这个两难消失。

第三样是分解失败本身成了信号。算法在某一步遇到非正的根号内量就停住，而这个停住是对输入矩阵的一次判决，不是算法对自身能力的抱怨。稍后会看到这句判决可以有三种截然不同的读法。

\begin{center}
\begin{tikzpicture}[x=1cm,y=1cm,font=\small]
  \node[font=\footnotesize,black!70] at (2,5.75) {一般方阵：\(PA=LU\)};
  \fill[orange!30] (0,5.2) -- (0,1.2) -- (4,1.2) -- cycle;
  \fill[blue!15]  (0,5.2) -- (4,5.2) -- (4,1.2) -- cycle;
  \draw[black!50] (0,1.2) rectangle (4,5.2);
  \draw[black!50] (0,5.2) -- (4,1.2);
  \node[font=\footnotesize] at (1.1,2.3) {\(L\)};
  \node[font=\footnotesize] at (2.9,4.1) {\(U\)};
  \node[font=\footnotesize,black!70] at (2,0.8) {存 \(n^2\) 个数，需选主元};
  \draw[fill=orange!45,draw=black!45] (0,0.05) rectangle (4,0.42);
  \node[right,font=\footnotesize,black!70] at (4.1,0.235) {\(\tfrac23n^3\)};
  \begin{scope}[shift={(6.9,0)}]
    \node[font=\footnotesize,black!70] at (2,5.75) {对称正定：\(A=LL^{\trans}\)};
    \fill[orange!30] (0,5.2) -- (0,1.2) -- (4,1.2) -- cycle;
    \draw[black!50] (0,1.2) rectangle (4,5.2);
    \draw[black!50] (0,5.2) -- (4,1.2);
    \node[font=\footnotesize] at (1.1,2.3) {\(L\)};
    \node[font=\footnotesize,black!55] at (2.9,4.25) {由对称性推知};
    \node[font=\footnotesize,black!55] at (2.9,3.8) {不算也不存};
    \node[font=\footnotesize,black!70] at (2,0.8) {存 \(n^2/2\) 个数，无需选主元};
    \draw[fill=orange!45,draw=black!45] (0,0.05) rectangle (2,0.42);
    \node[right,font=\footnotesize,black!70] at (2.1,0.235) {\(\tfrac13n^3\)};
  \end{scope}
\end{tikzpicture}

\emph{同一个 \(n\)，右边只算一个三角、只存一个三角，而且不必为选主元打乱行序。}
\end{center}

图右侧那半块空白比它看上去更值钱。省下的不只是存储与乘加，还有主元搜索带来的不规则访问；剩下的运算全是连续的、可分块的，实测的加速常常比 \(2:1\) 这个运算量之比更好。两千阶的核矩阵在这条路上是 \(2.7\times10^9\) 次浮点运算，一次三角求解只要 \(8\times10^6\) 次——几百轮超参数搜索能不能在一杯咖啡的时间里跑完，差别就在这里。

\subsection{根号下的数为什么一直是正的}

把 \(A=LL^{\trans}\) 逐项写开，\(a_{ij}=\sum_{k\le\min(i,j)}l_{ik}l_{jk}\)，按列解出来就是递推公式：

\[
\begin{aligned}
l_{jj}&=\Bigl(a_{jj}-\sum_{k<j}l_{jk}^2\Bigr)^{1/2},\\
l_{ij}&=\frac{a_{ij}-\sum_{k<j}l_{ik}l_{jk}}{l_{jj}},\qquad i>j.
\end{aligned}
\]

存在性的骨架只有两步。根号内那个量是消去前 \(j-1\) 列之后 Schur 补 \(S=A_{22}-A_{21}A_{11}^{-1}A_{12}\) 的首个对角元。对任意 \(x\ne0\) 取 \(y=-A_{11}^{-1}A_{12}x\)，把 \((y,x)\) 代进 \(A\) 的二次型展开，交叉项恰好抵掉，得

\[
x^{\trans}Sx=\begin{pmatrix}y\\x\end{pmatrix}^{\trans}A\begin{pmatrix}y\\x\end{pmatrix}>0 .
\]

于是 Schur 补仍正定，它的对角元 \(e_i^{\trans}Se_i\) 全为正，开方永远有意义；对列数归纳即得存在性，把对角元取正号就得到唯一性。

条件对账：论证用到 \(A_{11}\) 可逆，这由正定矩阵的主子阵仍正定保证；用到严格不等号，半正定只给出 \(\ge\)，根号内可以恰好为零，此时 \(L\) 存在但不唯一；还用到精确算术，浮点下根号内的量带着 \(O(nu\norm{A})\) 的误差，当它的真值本身小于这个量级时，正负就不再由数学决定，而由舍入决定。第三条正是第三十七轮那次异常的出处。

顺带落下一个副产品：\(\det A=\prod_i l_{ii}^2\)，取对数得 \(\log\det A=2\sum_i\log l_{ii}\)。

\subsection{一次分解，解、行列式与采样都在里面}

回到那个循环。有了 \(K=LL^{\trans}\)，边际似然要的两样东西都是现成的。\(K\alpha=y\) 走两次三角求解，各 \(n^2\) 次运算；\(\log\det K=2\sum_i\log l_{ii}\) 只需读一遍对角元再取对数——注意这条路避开了先算行列式再取对数的做法，两千个数连乘在双精度下必然下溢成零，取对数只能得到负无穷。预测方差需要的 \(k_*^{\trans}K^{-1}k_*\) 同样不必碰逆矩阵：解 \(Lv=k_*\)，答案就是 \(\norm{v}^2\)。

还有第四样，通常不在同一节课上讲。若 \(z\) 的各分量是独立标准正态，则

\[
\Cov(Lz)=\E\bigl[Lzz^{\trans}L^{\trans}\bigr]=LL^{\trans}=K .
\]

要从指定协方差的正态分布抽样，一次 Cholesky 加一次矩阵向量乘就够了。反过来看，\(L^{-1}(x-\mu)\) 把相关样本白化成独立标准正态，而 Mahalanobis 距离 \((x-\mu)^{\trans}K^{-1}(x-\mu)\) 正是白化后向量的平方范数。\hyperref[sec:05-Probability/05-Joint-Distributions/04]{``随机向量与协方差矩阵''}里的线性变换传播规则，到这里变成了两行可执行代码；同一个 \(L\) 既解方程又抽样，是因为它同时充当分解的因子和把标准正态拉成指定形状的那个变换。

\subsection{分解失败时该问的三个问题}

现在看第三十七轮。异常的字面意思只有一句：某一步根号内的量不为正。它背后至少有三类原因，处理方式完全不同。

第一类是矩阵本来就不该正定。样本协方差在样本数少于维数时必然奇异，这是秩的事实，加多少对角量都改不了它；\(A^{\trans}A\) 里某个下标写错、协方差用了有偏的估计器、核函数的参数落到了非法区间，都会产生同样的症状。这时该改的是构造矩阵的那段代码。

第二类是对称性在浮点里丢了。\(K\) 由 \(k(x_i,x_j)\) 逐项算出，理论上对称，实测的 \(\norm{K-K^{\trans}}/\norm{K}\) 却可能有 \(10^{-14}\) 量级；若中间做过 \(BCB^{\trans}\) 这类运算，偏差会更大。多数库只读下三角，于是这类误差被默默吞掉，但吞掉不等于无害——被读的那一半可能恰好是偏差较大的一半。先算一下这个比值，再决定要不要显式做 \(\tfrac12(K+K^{\trans})\)。

第三类是矩阵确实半正定或极接近。两个采样点重合或者相距远小于长度尺度时，\(K\) 的两行几乎相同，最小特征值掉到 \(u\norm{K}\) 以下，上一小节条件对账里的第三条就兑现了。只有这一类适合用加对角的办法处理，而且它不是一次数值补丁。

\subsection{jitter 改的是模型，不只是矩阵}

把 \(K\) 换成 \(K_\varepsilon=K+\varepsilon I\)，全部特征值抬高 \(\varepsilon\)，条件数从 \(\sigma_1/\sigma_n\) 降到 \((\sigma_1+\varepsilon)/(\sigma_n+\varepsilon)\)，取 \(\varepsilon\) 为 \(10^{-6}\norm{K}\) 量级通常就足以让分解通过。数值上这一步无可指摘。

统计上它写在一个已经有主人的位置。Gaussian process 回归的模型本来就是 \(K+\sigma_n^2I\)，那个对角项的身份是观测噪声的方差。所以以数值理由加进去的 \(\varepsilon\)，在似然里的身份是一个方差为 \(\varepsilon\) 的白噪声，不管加它的人是不是这么想的。后果有两条：\(\varepsilon\) 会进入预测方差，把不确定性整体抬高一点；而若超参数优化正在估计 \(\sigma_n^2\)，人为固定的 \(\varepsilon\) 相当于给它加了一个下界，估出来的噪声水平就不再是数据说的那个数。同样的结构在\hyperref[sec:13-Machine-Learning/09-Kernel-Methods-and-Gaussian-Processes/05]{``Gaussian process 把核变成函数先验''}里作为噪声项出现，在岭回归里作为正则化项出现，三者是同一个 \(+\lambda I\)。

因此 \(\varepsilon\) 的量级、加它的理由、以及结论对它是否敏感，都应当记录下来。把它当成不必写进方法一节的实现细节，等于在模型里塞进一个没有署名的假设。

\subsection{稀疏时成本不由维数决定，由消元顺序决定}

稠密情形的 \(\tfrac13n^3\) 是确定的。稀疏正定矩阵——有限元刚度阵、加了对角的图 Laplacian、稀疏精度矩阵——的成本却不由阶数和非零个数决定，而由填充决定：消去一个未知量，会在它的所有邻居之间连上新边，\(L\) 里因此冒出 \(A\) 中没有的非零元。同一个矩阵换一种未知量编号，\(L\) 的非零数可以差几个数量级。近似最小度、嵌套剖分这些排序算法，做的都是在消元图上找一个让填充尽量少的顺序；符号分解则先只跑图结构，算出 \(L\) 的非零模式与总代价，再决定值不值得做数值分解。

这里能这么做，还是靠前面第二样东西。一般 LU 的行序要同时服务于数值稳定与稀疏性，两个目标互相牵制，工业级求解器只好设阈值折中；对称正定不必选主元，排序就成了纯粹的组合问题，可以完全按稀疏性来挑。结构假设的红利在这里第二次兑现。

当填充大到连符号分解都劝退时，精确分解这条路就走到头了。丢掉一部分填充元，得到一个不完全的 \(\tilde L\)，用 \(\tilde L\tilde L^{\trans}\) 去近似 \(A\)，把它当预条件器交给只需要矩阵向量乘的迭代法——从这里开始，问题不再是怎样把矩阵拆开，而是怎样在根本不拆的前提下逼近解，那是\hyperref[sec:08-Numerical-Analysis/05-Numerical-Linear-Algebra/04]{``大型稀疏系统与 Krylov 迭代''}的主题。
